\section*{Chapitre \ref{ch:g2:addition} --- Addition avec retenue}

\begin{solution}{exo:g2:addition:1}
$34 + 25 = 59$~; \quad $126 + 43 = 169$~; \quad $507 + 82 = 589$.
\end{solution}

\begin{solution}{exo:g2:addition:2}
$47 + 8 = 55$~; \quad $56 + 27 = 83$~; \quad $65 + 19 = 84$.
\end{solution}

\begin{solution}{exo:g2:addition:3}
$268 + 145 = 413$~; \quad $374 + 158 = 532$~; \quad
$429 + 293 = 722$.
\end{solution}

\begin{solution}{exo:g2:addition:4}
$38 + 5 = 38 + 2 + 3 = 43$~; \quad $67 + 6 = 67 + 3 + 3 = 73$~; \quad
$54 + 8 = 54 + 6 + 2 = 62$.
\end{solution}

\begin{solution}{exo:g2:addition:5}
$42 + 36 = 70 + 8 = 78$~; \quad $53 + 24 = 70 + 7 = 77$~; \quad
$61 + 28 = 80 + 9 = 89$.
\end{solution}

\begin{solution}{exo:g2:addition:6}
$58 + 34$~: estimation $60 + 30 = 90$~; exactement $92$.
$296 + 187$~: estimation $300 + 200 = 500$~; exactement $483$.
\end{solution}

\begin{solution}{exo:g2:addition:7}
$128 + 95 = 223$. Ensemble, ils ont $223$ billes.
\end{solution}

\begin{solution}{exo:g2:addition:8}
Le premier nombre plus $27$ fait $64$, donc c'est
$64 - 27 = 37$~: le chiffre des unit\'es manquant est $7$ (le nombre
est $37$).
\end{solution}

\begin{solution}{exo:g2:addition:9}
Enfants~: $27 + 28 = 55$~; avec les adultes, $55 + 6 = 61$
personnes.
\end{solution}

\begin{solution}{exo:g2:addition:10}
$27$ et $33$ ont pour unit\'es $7$ et $3$ (amis de $10$)~:
$27 + 33 = 60$, puis $60 + 40 = 100$. Les trois nombres sont $27$,
$33$ et $40$.
\end{solution}
